\label{sec:signed306}

The main direct construction starts from the literal centered polynomial
of package 131~\cite{cap131}, in 324 unrestricted integer coordinates.  The source is
pinned to ProveIt commit
\srcpin; its equation, centering
certificate, and ordinary arithmetic proof are retained in the archive.
The new polynomial is
\begin{equation}\label{eq:306definition}
 P_{306}(z_0,\ldots,z_{305})=\sum_{j=0}^{205}r_j(z)^2,
 \qquad z\in\mathbb Z^{306}.
\end{equation}
Every $r_j$ is an explicitly listed integer polynomial of degree at most two.  There
are no implicit domain predicates: all 306 coordinates are unrestricted.
The literal residuals and the complete collected coefficient list are
supplied, together with their independent verifier.

\begin{center}
\begin{tabular}{lrrrrr}
\toprule
Equation & Variables & Rows & Degree & Terms & Height\\
\midrule
Pinned parent &324&208&4&5808&128446\\
New equation &306&206&4&6259&128446\\
\bottomrule
\end{tabular}
\end{center}

Thus 18 coordinates and two rows are removed at the cost of 451 collected
terms; both coefficient height and residual height, respectively 128446
and 256, are preserved.  The comparison is within this direct C222
construction.  Much smaller variable counts exist in the repository with
different degrees and coefficient architectures.

\begin{theorem}\label{thm:306-equivalence}
Let $P_{324}$ be the pinned parent polynomial.  Ordinary mathematics
formalizable in ZFC proves
\[
 (\exists z\in\mathbb Z^{306}\;P_{306}(z)=0)
 \quad\Longleftrightarrow\quad
 (\exists w\in\mathbb Z^{324}\;P_{324}(w)=0).
\]
Consequently, composing with the retained parent calibration,
\[
 \mathrm{ZFC}\vdash
 \bigl[\neg\exists z\in\mathbb Z^{306}\;P_{306}(z)=0\bigr]
 \longleftrightarrow\operatorname{Con}(\mathrm{ZFC}).
\]
\end{theorem}

The new idea is a linear orientation term inside the existing bounded
Pell-index congruence.  It permits nine sign graphs to be replaced by one
strict product.  The proof must recover the signs before invoking the
parent's complete arithmetic theorem.  We give that order explicitly.

\subsection{Linear orientation in a bounded residue pool}

For $A\geq2$ let $X_j(A),Y_j(A)$ be the canonical Pell sequences:
\[
 X_0=1,\quad Y_0=0,\qquad
 X_{j+1}=AX_j+(A^2-1)Y_j,\quad Y_{j+1}=X_j+AY_j.
\]
The unchanged parent square frame and complete anchor yield, independently
of the nine encoded-coordinate signs,
\begin{equation}\label{eq:306frame}
 W\geq258,\qquad A>W^W,\qquad A-1>4W^3.
\end{equation}
They also give 15 principal Pell norms, absolute indices $0\leq j_i<W$,
signed proposed indices $|k_i|<W$, and absolute bounds $|n|,|m|<W$ for
every requested base and output.  The fully indexed anchor is at $W$;
its unchanged cap bounds every principal ordinate in absolute value.
Its proof uses the frame, the anchor common-multiplier graph and the cap
graph, all retained here.  No $k_i\geq0$ premise is required.

Pell classification supplies signs $\sigma_i,\epsilon_i\in\{1,-1\}$ with
\[
 x_i=\sigma_iX_{j_i}(A),\qquad y_i=\epsilon_iY_{j_i}(A).
\]
At $j_i=0$ the choice of $\epsilon_i$ is immaterial.  Modulo $A-1$,
the recurrence gives $X_j\equiv1$ and $Y_j\equiv j$.

\begin{lemma}[Linear orientation pool]\label{lem:306pool}
Assume the preceding bounds.  Let $I$ be a set of nine of the fifteen
principal pairs.  The single row
\begin{equation}\label{eq:306pool}
 \sum_{i=1}^{15}(y_i-k_i)^2+\sum_{i\in I}(1-x_i)=(A-1)q_0,
 \qquad q_0\in\mathbb Z,
\end{equation}
forces $\epsilon_i j_i=k_i$ for every $i$ and $\sigma_i=1$ for $i\in I$.
\end{lemma}

\begin{proof}
Reduce \eqref{eq:306pool} modulo $A-1$.  It makes $A-1$ divide
\[
 E=\sum_{i=1}^{15}(\epsilon_i j_i-k_i)^2+
     \sum_{i\in I}(1-\sigma_i).
\]
Every term in the last sum belongs to $\{0,2\}$.  Therefore
\[
 0\leq E\leq60(W-1)^2+18<4W^3<A-1.
\]
For example,
$4W^3-60(W-1)^2-18=4W^2(W-15)+120W-78>0$ for $W\geq258$.
Hence $E=0$, and every nonnegative summand vanishes.
\end{proof}

Only the residue representative $E$ is small and nonnegative.  The actual
left side of \eqref{eq:306pool} need not be either.  In particular the pool
quotient must remain unrestricted.  A squared term $(x_i-1)^2$ was used
in the parent's conservative orientation variant; the linear term
$1-x_i$ works because its residue has the nonnegative alphabet $\{0,2\}$.
The new application combines this cheaper orientation with removal of
the nine encoded-coordinate sign graphs.

\subsection{Signed indices and the early bootstrap}

We need a power lemma that treats zero and negative proposed indices.
Write $Y_{-j}=-Y_j$ and $X_{-j}=X_j$.

\begin{lemma}[Signed-index power request]\label{lem:306request}
Let $2\leq n<W$, $|k|<W$, $|m|<W$, and put
$T=2An-n^2-1$.  Suppose
\[
 y=Y_k(A),\qquad x=\pm X_{|k|}(A),\qquad
 x+(n-A)y=m+Tz
\]
for an integer $z$.  If $x=+X_{|k|}$, then $k\geq0$ and $m=n^k$.
If $k>0$, either choice initially allowed for $x$ is forced to be positive.
At $k=0$ and positive $x$, the output is exactly $m=1$.
\end{lemma}

\begin{proof}
For $j\geq0$ put
$H_j=X_j+(n-A)Y_j$ and $J_j=X_j-(n-A)Y_j$.  Direct multiplication and
the recurrence give
\[
 H_jJ_j=1+TY_j^2,\qquad H_{j+1}=nH_j+TY_j,
\]
so $H_j\equiv n^j\pmod T$.  The frame bounds imply $T>W^W$.  For
$1\leq j<W$ the required differences have absolute value below $W^W$;
in particular
$|mn^j-1|\leq(W-1)^W+1<W^W$.

If $k=-j<0$ and $x=+X_j$, the request gives
$mn^j\equiv1\pmod T$, hence $mn^j=1$ as an integer equality.  This is
impossible since $n^j\geq2$.  If $k=j>0$ and $x=-X_j$, it similarly gives
$mn^j=-1$, also impossible.  The remaining positive-index branch gives
$m\equiv n^k\pmod T$, and the size bound makes it equality.  At $k=0$,
$x=1,y=0$ gives $m=1$ directly by the same small-residue argument.
\end{proof}

The selected nine pairs are precisely the requests at base $q+1$ for
the eight encoded CRT columns and their reciprocal-index column.  The
following bootstrap does not use any of their removed sign graphs.
The retained rows and retained nonnegative quantities $r,s,o$ give
\[
 t=16+h_0^2\geq16,\qquad B=r+s+1\geq1,\qquad q=B+o+1\geq2.
\]
The t-pair has its signed intended index from Lemma~\ref{lem:306pool}.
Its genuine request at base $q$ and positive index $t$ forces its abscissa
positive by Lemma~\ref{lem:306request}.  The same pair therefore gives
\[
 V_0=(F^2+1)^t,
\]
including the possible base-one case before the sign of $F$ is known.
For that case the same recurrence gives $H_t\equiv1\pmod{2(A-1)}$;
the bound $|V_0|<W$ then forces $V_0=1$.
The literal ceiling rows give $FB=V_0+s>0$, so $F\geq1$.  With
$b=F^2+1$, both $b>F$ and $b^{t-1}>F$, whence
\begin{equation}\label{eq:306Bbound}
 FB\geq b^t>Fb,\qquad B>F^2+1.
\end{equation}
The unchanged cofinal row $y=V_0(V_0-1)$ also gives $y>F$; the x-sign
graph remains present.  Neither factorial meaning nor CRT interpretation
has been used at this stage.

Lemma~\ref{lem:306pool} orients the nine selected abscissas, and
Lemma~\ref{lem:306request}, at base $q+1\geq3$, gives
\[
 k_a,k_b,k_c,k_d,k_S,k_{S'},k_u,k_{v_q},k_{\mathrm{index}}\geq0.
\]
Zero remains a genuine possibility at this point.  The reciprocal-index
rows give
\[
 B\,v=F^2+k_{\mathrm{index}}+1
\]
for an unrestricted integer $v$.  If $k_{\mathrm{index}}=0$, its right
side lies strictly between zero and $B$ by \eqref{eq:306Bbound}, an
impossible integer multiple of $B$.  Therefore $k_{\mathrm{index}}\geq1$.

\subsection{Six products and one fused positivity row}

The other eight indices are nonnegative.  Introduce six unrestricted
coordinates $P_2,\ldots,P_7$ with
\begin{align*}
 P_2&=k_a k_b,&P_3&=P_2 k_c,&P_4&=P_3 k_d,\\
 P_5&=P_4 k_S,&P_6&=P_5 k_{S'},&P_7&=P_6 k_u.
\end{align*}
Finally introduce three unrestricted integers $a,b,c$ and impose
\begin{equation}\label{eq:306product}
 P_7 k_{v_q}-1=a^2+b^2+c(c+1).
\end{equation}
The right side is nonnegative at every integer triple.  The product of
the eight nonnegative indices is thus at least one, so every index is
at least one.  This restores all nine old encoded-coordinate signs.
Every displayed row is quadratic; a separate $P_8$ coordinate is avoided.

Conversely every nonnegative integer $n$ is
$a^2+b^2+c(c+1)$.  Indeed, $4n+1$ is not of the forbidden form
$4^u(8v+7)$, so Legendre's three-square theorem represents it by three
squares.  Exactly one is odd; writing the three terms as
$(2a)^2,(2b)^2,(2c+1)^2$ gives the result.  Thus every old zero, whose
eight indices are positive, has witnesses for \eqref{eq:306product}.

The eight-factor idea has a predecessor in the repository's
\texttt{signed-pell-cap-refined-scalar-product} proof.  That proof obtains
nonnegative indices from a different residue-one chart.  Here the original
raw remainder chart is retained; the linear orientation pool supplies the
required signs directly.  The change is this complete composition and its
literal smaller polynomial, not invention of positive-product reasoning.

\subsection{Literal chart and exact witness transformations}

The centered parent has the following requested-index chart.  The third
column lists its nine selected abscissa coordinates; all have zero
centering shift.
\begin{center}
\begin{tabular}{lrr}
\toprule
Index & Expression & Abscissa coordinate\\
\midrule
$k_a$ &$z_{10}-1$&174\\
$k_b$ &$z_{11}$&180\\
$k_c$ &$z_{12}$&184\\
$k_d$ &$z_{13}$&188\\
$k_S$ &$z_{14}$&192\\
$k_{S'}$ &$z_{15}$&200\\
$k_u$ &$z_{16}$&204\\
$k_{v_q}$ &$z_{17}$&212\\
$k_{\mathrm{index}}$ &$z_{18}$&216\\
\bottomrule
\end{tabular}
\end{center}
The old nonnegative quantities are $k_i-1$.  Delete precisely parent rows
184--192, the nine scalar graphs labeled \texttt{domain\_guard.3} through
\texttt{domain\_guard.11}, and their 27 private coordinates 252--278.
They occur nowhere else.  Modify parent row 180, the pooled-index row, by
adding precisely
\[
 9-z_{174}-z_{180}-z_{184}-z_{188}-z_{192}
   -z_{200}-z_{204}-z_{212}-z_{216}.
\]
The private pool quotient is coordinate 242 and occurs only in this row.
The actual Pell center is $A=z_{151}+1$, so its divisor is $A-1=z_{151}$.

For a new zero, the preceding argument proves the nine missing signs.
Let $T_0=\sum_{i\in I}(1-X_i)$.  Each selected abscissa is canonical,
so $T_0/(A-1)$ is an integer.  Set
\[
 q_{\mathrm{old}}=q_{\mathrm{new}}-T_0/(A-1),
\]
and choose fresh Legendre triples for all nine deleted scalar graphs.
Every other retained parent coordinate stays fixed.  The private quotient
change restores the old pooled row and affects no other row.  We now have
a literal zero of the entire parent polynomial, so its complete C222
calibration may be invoked.

For the reverse direction, a parent zero has positive intended indices
and canonical abscissas by its ordinary power theorem.  Set
\[
 q_{\mathrm{new}}=q_{\mathrm{old}}+T_0/(A-1),
\]
forget the 27 old guard witnesses, calculate the six product coordinates,
and supply a triple for \eqref{eq:306product}.  All new rows hold.  These
maps preserve the native front and every retained parent coordinate except
the private pool quotient.  The Legendre witnesses are not unique, so
there is no claim of a bijection between individual full zero tuples.
This proves Theorem~\ref{thm:306-equivalence}.

The parent projection and centering are also part of the supplied
certificate boundary.  The nine deleted guard rows survived projection;
no restoration polynomial depends on their removed auxiliary coordinates
or the private pool quotient.  The independently checked occurrence sets
therefore justify using the parent frame and anchor before restoring those
graphs, without a circular appeal to the entire parent theorem.

\subsection{Finite certification and its scope}

There are $324-27+6+3=306$ coordinates and $208-9+6+1=206$ residuals.
Their sum of squares has a nonzero degree-four coefficient, so the degree
is exactly four.  The independent verifier imports no producer code.  It
checks every retained row, the complete changed pool row, every removed
witness occurrence set, the quotient's privacy, the affine chart, the
six products, the fused positivity row, all collected coefficients and
the exact metrics.  It also rejects deliberate changes to an orientation
term, a factor shift, quotient privacy, and an expanded coefficient.

These finite checks certify literal algebra and artifact bindings.  They
do not establish unbounded number theory by sampling.  Pell classification,
the retained full anchor, Legendre's theorem and the ordinary parent
C222/ZF/ZFC calibration remain the identified mathematical dependencies.
No new end-to-end proof-assistant certificate or encoded object-ZFC proof
is claimed.
